\chapter{Analysis / Software Requirements Specification (SRS)}\doublespacing

\section{Introduction to SRS}
The Software Requirements Specification (SRS) establishes a formal, rigorous blueprint for the design, architecture, implementation, and verification of the \textbf{LightAngo} compiler system. It details the complete functional expectations, non-functional performance and reliability criteria, target hardware and software operational environments, and detailed use case scenarios governing system interactions.

This specification serves as the fundamental contract between language design specifications, engineering constraints, and verification metrics for the ahead-of-time (AOT) compilation of the LightAngo programming language.

\section{System Overview}
The LightAngo system is a self-contained, native compiler engineered in C++20 for 64-bit Linux architectures. It processes source files written in the custom LightAngo high-level programming language (`.lango` or `.ango` files) and translates them through a multi-stage compilation pipeline into 64-bit x86 NASM assembly (`.asm`), assembling and linking them into executable ELF binaries.

Unlike conventional compilers that treat comments purely as whitespace to be ignored, LightAngo integrates an active \textit{Documentation Verification Engine} directly into its recursive-descent parsing and semantic analysis layers. The system evaluates both syntactic correctness and documentation compliance before permitting assembly synthesis.

\section{Functional Requirements}
The functional requirements specify the exact operations, transformations, and behaviors the LightAngo compiler must execute across each phase of the compilation lifecycle.

\begin{enumerate}
    \item \textbf{FR-01: Source Code Ingestion}
    \begin{itemize}
        \item The system shall accept standard UTF-8 encoded text source files containing valid or invalid LightAngo source code.
        \item The system shall support command-line arguments specifying input files, output targets, and verbose compilation flags.
    \end{itemize}

    \item \textbf{FR-02: Lexical Analysis and Tokenization}
    \begin{itemize}
        \item The lexical scanner shall convert the raw character stream into strongly typed tokens, including keywords (`let`, `if`, `elif`, `else`, `while`, `for`, `fn`, `exit`), identifiers, integer/string literals, arithmetic operators (`+`, `-`, `*`, `/`, `\%`), comparison operators (`==`, `!=`, `<`, `<=`, `>`, `>=`), and punctuation (`(`, `)`, `{`, `}`, `;`, `,`).
        \item The scanner shall track precise source coordinates (line number and column offset) for every token to support pinpoint error reporting.
        \item The scanner shall capture single-line comments (`//...`) and multi-line comments (`/* ... */`) into a structured comment buffer linked to line numbers rather than discarding them.
    \end{itemize}

    \item \textbf{FR-03: Syntax Analysis and AST Generation}
    \begin{itemize}
        \item The recursive descent parser shall validate the token stream against the formal Context-Free Grammar (CFG) of the LightAngo language.
        \item The parser shall construct a strongly-typed Abstract Syntax Tree (AST) representing statements, binary/unary expressions, variable bindings, conditional branches, and function definitions.
        \item The parser shall enforce operator precedence and associativity rules (multiplication/division over addition/subtraction).
    \end{itemize}

    \item \textbf{FR-04: Grammar-Level Comment Enforcement}
    \begin{itemize}
        \item Prior to parsing any loop construct (`while`, `for`) or function declaration (`fn`), the parser shall inspect the preceding comment registry.
        \item If no valid, non-empty documentation comment is detected within the immediate preceding source lines, the compiler shall generate a fatal diagnostic error (e.g., \texttt{"Error: Missing mandatory documentation comment before loop construct at line X"}).
        \item Compilation shall be immediately gated, preventing binary emission for undocumented code.
    \end{itemize}

    \item \textbf{FR-05: Semantic Analysis and Scoped Symbol Table Management}
    \begin{itemize}
        \item The semantic analyzer shall maintain a hierarchical, scoped symbol table tracking variable identifiers, data types, stack memory offsets, and mutability flags.
        \item The system shall detect and reject undeclared identifier references, duplicate declarations in the same lexical scope, and invalid type assignments.
        \item The system shall verify control flow integrity, ensuring that system calls and return pathways conform to language semantics.
    \end{itemize}

    \item \textbf{FR-06: Machine-Independent Optimization}
    \begin{itemize}
        \item The compiler shall perform constant folding on compile-time evaluable expressions (e.g., evaluating \texttt{(10 - 2 * 3) / 2} to \texttt{2} at compile time).
        \item The system shall eliminate redundant stack pushes and pops for trivial sequential expressions.
    \end{itemize}

    \item \textbf{FR-07: x86-64 NASM Code Generation}
    \begin{itemize}
        \item The code generation engine shall traverse the validated AST and emit valid 64-bit Netwide Assembler (NASM) code conforming to the Linux System V AMD64 ABI.
        \item The generator shall manage local variable allocation on the stack frame relative to the `%rsp` register.
        \item The generator shall emit correct conditional jump labels (`je`, `jne`, `jl`, `jge`, `jmp`) for `if-elif-else` constructs and loop iterations.
        \item The generator shall correctly synthesize Linux system call sequences (e.g., `mov rax, 60` and `mov rdi, [val]; syscall` for process exit).
    \end{itemize}

    \item \textbf{FR-08: Assembler and Linker Integration}
    \begin{itemize}
        \item The system shall automatically invoke `nasm -f elf64 out.asm -o out.o` to produce 64-bit object files.
        \item The system shall automatically invoke GNU Linker `ld -o out out.o` to produce executable ELF binaries.
    \end{itemize}

    \item \textbf{FR-09: Contextual Error Reporting and Diagnostics}
    \begin{itemize}
        \item The diagnostic engine shall print descriptive, color-coded error messages containing the exact file name, line number, column number, offending token, and an explanatory remediation hint.
    \end{itemize}
\end{enumerate}

\section{Non-Functional Requirements}
Non-functional requirements specify the quality attributes, operational limits, and performance benchmarks of the compiler system.

\begin{enumerate}
    \item \textbf{Performance and Throughput:}
    \begin{itemize}
        \item The compiler shall process at least 50,000 lines of source code per second on standard desktop hardware.
        \item AST generation and semantic analysis for typical programs ($<1000$ lines) shall complete in under 50 milliseconds.
    \end{itemize}

    \item \textbf{Reliability and Determinism:}
    \begin{itemize}
        \item The compiler shall produce strictly deterministic machine code across repeated compilation runs on identical source inputs.
        \item The compiler shall never terminate abruptly due to unhandled exceptions, memory segmentation faults, or buffer overflows on invalid user inputs.
    \end{itemize}

    \item \textbf{Memory Safety and Resource Efficiency:}
    \begin{itemize}
        \item All internal AST structures and symbol table allocations shall utilize modern C++ smart pointers (`std::unique_ptr`, `std::shared_ptr`) or arena allocators to ensure zero memory leaks.
        \item Peak memory consumption during compilation shall remain below 64 MB for standard input files.
    \end{itemize}

    \item \textbf{Usability and Developer Feedback:}
    \begin{itemize}
        \item Command-line interface (CLI) shall be intuitive, providing informative help flags (`--help`, `--version`, `--verbose`).
        \item Error diagnostics shall follow standardized compiler formatting compatible with modern IDEs (CLion, VS Code, Vim).
    \end{itemize}

    \item \textbf{Portability and System Compatibility:}
    \begin{itemize}
        \item The compiler shall build cleanly using GCC (g++ 11+) or Clang (clang++ 14+) with standard CMake build scripts.
        \item Generated assembly shall run natively on any x86-64 Linux kernel (v4.x, v5.x, v6.x).
    \end{itemize}
\end{enumerate}

\section{Hardware Requirements}
The minimum and recommended hardware specifications for building and running the LightAngo compiler are summarized in Table~\ref{tab:hardware_req}.

\begin{table}[h!]
\centering
\caption{Hardware Requirements} \label{tab:hardware_req}
\setstretch{1.3}
\begin{tabular}{|l|l|l|}
\hline
\textbf{Component} & \textbf{Minimum Specification} & \textbf{Recommended Specification} \\
\hline
Processor & Intel Core i3 / AMD Ryzen 3 (64-bit) & Intel Core i5/i7 / AMD Ryzen 5/7 \\
Clock Speed & 2.0 GHz Dual-Core & 3.0 GHz Quad-Core or higher \\
RAM & 4 GB DDR4 & 8 GB / 16 GB DDR4/DDR5 \\
Free Storage & 5 GB free disk space & 20 GB SSD storage \\
Architecture & x86-64 (AMD64) & x86-64 with AVX2 support \\
\hline
\end{tabular}
\end{table}

\section{Software Requirements}
The software dependencies and development toolchain specifications are detailed in Table~\ref{tab:software_req}.

\begin{table}[h!]
\centering
\caption{Software Requirements} \label{tab:software_req}
\setstretch{1.3}
\begin{tabular}{|l|l|}
\hline
\textbf{Category} & \textbf{Tool / Environment Specification} \\
\hline
Operating System & Linux (Ubuntu 20.04/22.04 LTS, Debian 11/12, or Fedora 38+) \\
Compiler Implementation Language & C++ (ISO C++20 standard) \\
Host Compiler Toolchain & GCC (g++ 11.0+) / Clang (clang++ 14.0+) \\
Integrated Development Environment & JetBrains CLion / Visual Studio Code \\
Build Automation & CMake (version 3.20 or higher) \\
Assembler & Netwide Assembler (NASM version 2.15+) \\
Linker & GNU Linker (\texttt{ld} from GNU binutils) \\
Version Control & Git \& GitHub \\
\hline
\end{tabular}
\end{table}

\section{Use Case Descriptions}
The operational workflows of the compiler are structured around nine core use cases spanning user interaction and compiler subsystem execution:

\begin{table}[h!]
\centering
\caption{Use Case Specifications} \label{tab:use_cases}
\setstretch{1.2}
\begin{tabular}{|p{2.0cm}|p{3.5cm}|p{8.5cm}|}
\hline
\textbf{Use Case ID} & \textbf{Use Case Name} & \textbf{Description} \\
\hline
\textbf{UC-01} & Write Source Code & Developer writes program statements, expressions, loops, and comments using LightAngo grammar. \\ \hline
\textbf{UC-02} & Ingest Source File & Compiler reads file from disk, loads content into memory buffer, and initializes line trackers. \\ \hline
\textbf{UC-03} & Lexical Scanning & Scanner converts character buffer into token list and captures comment structures. \\ \hline
\textbf{UC-04} & Recursive Parsing & Parser validates token sequences against grammar rules and builds AST hierarchy. \\ \hline
\textbf{UC-05} & Comment Validation & Verification engine inspects presence of comments before loops and function nodes. \\ \hline
\textbf{UC-06} & Semantic Validation & Analyzer resolves scopes, checks type rules, and verifies variable bindings in symbol table. \\ \hline
\textbf{UC-07} & Assembly Emission & Generator synthesizes x86-64 NASM instructions from validated AST. \\ \hline
\textbf{UC-08} & Binary Linkage & Toolchain invokes NASM assembler and GNU linker to produce final ELF binary executable. \\ \hline
\textbf{UC-09} & Diagnostic Output & Compiler displays structured error messages or execution status feedback to the user. \\ \hline
\end{tabular}
\end{table}

\newpage
